\documentclass[11pt,a4paper]{article}

\usepackage[top=2.5cm,bottom=2.5cm,left=2.5cm,right=2.5cm]{geometry}
\usepackage{times}
\usepackage{booktabs}
\usepackage{amsmath}
\usepackage{hyperref}
\usepackage{url}
\usepackage{microtype}
\usepackage{caption}
\usepackage{parskip}

\hypersetup{
  colorlinks=true,
  linkcolor=black,
  citecolor=black,
  urlcolor=black
}

\title{\textbf{Experience Breeds Skepticism: How Professional\\
Coding Experience Predicts Distrust in AI-Generated\\
Output Accuracy Among Software Developers}}

\author{Anonymous Author\\
\texttt{anonymous@institution.edu}}

\date{}

\begin{document}

\maketitle
\thispagestyle{empty}

%-----------------------------------------------------------------------
\begin{abstract}
As AI-powered code generation tools become ubiquitous in software development,
understanding which developers trust their outputs is critical for adoption
strategy and risk management. A common assumption holds that technically
experienced developers will trust AI tools more because they have greater
capacity to evaluate and integrate AI suggestions effectively. Using data from
the \emph{Stack Overflow Developer Survey 2024} ($N = 65{,}437$), this study
examines the relationship between professional coding experience and self-reported
trust in AI output accuracy among active AI tool users ($n = 31{,}025$). Contrary
to the intuitive expectation, a significant positive association between experience
and distrust emerges: developers in the \emph{Highly distrust} category average
$10.25$ years of professional experience versus $8.75$ years in the most trusting
group. A Kruskal--Wallis test confirms the difference across five trust levels
($H(4) = 127.78$, $p < 0.001$), and Spearman rank correlation yields
$r_s = 0.10$ ($p < 0.001$). Quartile analysis shows that novice developers
($0$--$3$ years) are disproportionately trusting ($53.1\%$ net trust against
$24.2\%$ net distrust), whereas senior developers (${\geq}\,13$ years) exhibit
near parity ($38.2\%$ vs.\ $35.2\%$). Education level does not account for this
pattern. These findings suggest that expertise-induced calibration --- not
technological enthusiasm --- governs how experienced practitioners assess
AI tool reliability.
\end{abstract}

%-----------------------------------------------------------------------
\section{Introduction}
\label{sec:intro}

The adoption of large-language-model (LLM)-based coding assistants has accelerated
dramatically since the public release of GitHub Copilot in 2022 and the subsequent
proliferation of tools such as ChatGPT, Amazon CodeWhisperer, and Cursor. According
to the \emph{Stack Overflow Developer Survey 2024} \cite{stackoverflow2024}, more than
$57\%$ of the $65{,}437$ respondents currently incorporate AI tools into their
development workflow --- a majority that crosses professional, regional, and
organisational lines. This rapid adoption has outpaced the research community's
understanding of how developers relate to the outputs these tools produce.

Trust in AI output accuracy is a critical moderator of tool effectiveness. Over-trusting
developers risk shipping incorrect AI-generated code without adequate review, introducing
latent bugs or security vulnerabilities. Under-trusting developers forgo productivity
benefits by redundantly verifying every suggestion, negating the promised efficiency gains
\cite{cai2019hello}. The distribution of trust across the developer population is therefore
a question of both theoretical and practical importance: if trust is unevenly distributed
in systematic ways, organisations can tailor onboarding, training, and code-review policies
to match the risk profile of their team composition.

Prior research on AI-assisted programming has focused mainly on productivity metrics and
usability outcomes \cite{vaithilingam2022expectation, ziegler2022productivity, bird2022copilot}.
These studies find that developers accept a substantial fraction of AI suggestions outright,
raising questions about the quality of human oversight. Less studied is the question of
\emph{who} trusts AI and \emph{why}. The present study addresses this gap by examining
one of the most fundamental individual-difference variables in software engineering:
professional coding experience.

A naive expectation might be that more experienced developers trust AI tools more. They
understand the underlying technology better, they have more context for evaluating whether
a suggestion is plausible, and they are more capable of integrating AI outputs into
existing codebases. An alternative --- and we argue more nuanced --- hypothesis is that
experience leads to \emph{better calibration}: experienced developers have seen enough
correct and incorrect code to detect AI errors that novices miss. This expertise-induced
skepticism would manifest as higher distrust among senior practitioners, even as they
continue to use the tools \cite{liao2023transparency}.

The \emph{Stack Overflow Developer Survey 2024} provides a uniquely suitable dataset to
adjudicate between these hypotheses. With over 65,000 responses globally, rich demographic
and attitudinal data, and direct questions about AI tool adoption and accuracy trust, it
enables large-scale observational analysis that no single-institution study could replicate.

\subsection{Research Question}

\emph{Does professional coding experience predict skepticism toward AI-generated output
accuracy among developers who currently use AI tools?}

This question is non-trivial because competing theoretical arguments generate opposite
predictions. The ``competence-based trust'' hypothesis (experienced developers understand
AI better $\rightarrow$ trust more) conflicts with the ``calibration'' hypothesis
(experienced developers detect AI errors better $\rightarrow$ trust less). The present
study tests these empirically.

\subsection{Hypotheses}

\noindent\textbf{H\textsubscript{0}:} Professional coding experience does not differ
across levels of self-reported trust in AI output accuracy. The mean years of experience
are equal across all five trust levels ($\mu_1 = \mu_2 = \mu_3 = \mu_4 = \mu_5$).

\noindent\textbf{H\textsubscript{1}:} Developers with more professional experience
exhibit systematically greater distrust of AI output accuracy. Mean professional
experience increases monotonically from the most trusting to the most distrusting group
($\mu_1 \leq \mu_2 \leq \mu_3 \leq \mu_4 \leq \mu_5$, with at least one strict inequality).

\subsection{Related Work}

Vaithilingam et al.\ \cite{vaithilingam2022expectation} conducted a controlled study
in which software developers used Copilot to complete programming tasks. They found that
participants frequently accepted syntactically correct but semantically flawed suggestions,
particularly when working in unfamiliar problem domains. Critically, those with relevant
expertise in the domain were faster to identify and reject incorrect suggestions ---
providing a plausible mechanism for experience-driven distrust at scale.

Ziegler et al.\ \cite{ziegler2022productivity} analysed acceptance rates of GitHub Copilot
suggestions across a large cohort of professional developers and found that acceptance
rates declined as developers' proficiency in the relevant programming language increased.
Although the primary outcome was productivity rather than trust, the acceptance-rate
gradient is consistent with the calibration hypothesis: knowing a language well enough
means knowing when an AI suggestion is wrong.

Bird et al.\ \cite{bird2022copilot} documented early Copilot adoption patterns among
engineers at Microsoft. Senior engineers with longer tenure tended to use AI suggestions
more selectively, often discarding proposals and re-prompting rather than accepting them
outright, while junior engineers showed higher raw acceptance rates. This behavioural
pattern aligns with the attitudinal differences the present study measures.

On the AI transparency side, Liao and Vaughan \cite{liao2023transparency} argue that
``appropriate reliance'' on AI systems requires users to have accurate mental models of
system capabilities and failure modes. Developing such models requires sustained experience
with the system --- but that experience initially manifests as \emph{increased} distrust
before possibly stabilising or improving as users learn which task types AI handles well.

Finally, the calibration literature in human--AI teaming consistently finds that domain
expertise enables practitioners to detect AI errors that novices miss, resulting in reduced
over-reliance over time \cite{cai2019hello}. Cai et al.\ demonstrate this in a medical
context, where experienced clinicians showed better-calibrated trust in AI diagnostic
recommendations than trainees. Whether this expertise-calibration relationship generalises
to software development has not been tested at scale --- the present study provides the
first large-scale evidence.

%-----------------------------------------------------------------------
\section{Methodology}
\label{sec:method}

\subsection{Dataset}

The analysis uses the publicly available \emph{Stack Overflow Developer Survey 2024}
\cite{stackoverflow2024}, an annual worldwide survey of software practitioners conducted
online between May and June 2024. The public release contains $65{,}437$ valid responses
from participants in over 185 countries. Respondents completed the survey voluntarily and
without financial compensation. Sampling is non-probabilistic (opt-in via the Stack Overflow
platform and social media), which has implications for generalisability discussed in
Section~\ref{sec:discussion}.

The survey instrument covers a wide range of topics including professional background,
technology and tool usage, compensation, and --- prominently in the 2024 edition ---
attitudes toward AI development tools. The AI attitude battery was newly expanded in 2024
to capture adoption, sentiment, perceived benefit, accuracy trust, and job threat perception,
making this wave particularly suitable for the present analysis.

\subsection{Variables}

\textbf{Dependent variable --- AI Accuracy Trust} (\texttt{AIAcc}): A five-point
ordinal survey item: ``How much do you trust the accuracy of the output from AI tools
as part of your development workflow?'' Response options were \emph{Highly trust},
\emph{Somewhat trust}, \emph{Neither trust nor distrust}, \emph{Somewhat distrust},
and \emph{Highly distrust}. This item was presented only to respondents who reported
currently using AI tools (\texttt{AISelect = Yes}). Respondents who reported not using
AI tools --- either planning to or not planning to --- did not receive this question,
which is why the \texttt{AIAcc} variable contains approximately $28{,}000$ missing values
in the public dataset.

\textbf{Independent variable --- Professional Coding Experience} (\texttt{YearsCodePro}):
The survey question reads: ``NOT including education, how many years have you coded
professionally (as part of your work)?'' Responses are binned: ``Less than 1~year'',
followed by integer years from 1 to 50, and ``More than 50 years''. For analysis,
``Less than 1 year'' was coded as $0$. The category ``More than 50 years'' ($n < 50$)
was excluded because it does not admit a precise numeric value and constitutes an extreme
tail with negligible influence on group-level statistics. The resulting distribution is
strongly right-skewed: mean $= 9.12$ years, SD $= 8.33$, median $= 6$ years, with
$37.1\%$ of the analytic sample reporting ten or more years of professional experience.

\textbf{Control variable --- Education Level} (\texttt{EdLevel}): An ordinal item with
seven categories from \emph{Primary/elementary school} (coded~0) to \emph{Professional
degree including Ph.D.} (coded~6). Education level was examined as a potential confounder
because it correlates with both career tenure and technology literacy.

\subsection{Analytic Sample}

The analysis is restricted to respondents with non-missing values on both \texttt{AIAcc}
and \texttt{YearsCodePro}, yielding an analytic sample of $n = 31{,}025$. All such
respondents also reported currently using AI tools (\texttt{AISelect = Yes}), by the
survey design described above. Table~\ref{tab:main} reports sample sizes per trust group,
ranging from $n = 699$ (\emph{Highly trust}) to $n = 12{,}175$ (\emph{Somewhat trust}).

\subsection{Analytic Strategy}

Because \texttt{YearsCodePro} is strongly right-skewed and \texttt{AIAcc} is ordinal,
the standard parametric one-way ANOVA assumptions are violated. Accordingly, the primary
inferential test is the non-parametric \textbf{Kruskal--Wallis H-test}, which tests
whether the rank distribution of professional experience differs across the five trust
levels. Effect size is reported as eta-squared ($\eta^2 = \text{SS}_\text{between} /
\text{SS}_\text{total}$) and Cohen's $d$ for key pairwise comparisons. Ties are handled
by averaging ranks.

A \textbf{Spearman rank correlation} $r_s$ quantifies the monotonic association between
trust ordinal (coded $1 = \text{Highly trust}$ through $5 = \text{Highly distrust}$)
and professional experience, assessed via the $t$-approximation
$t = r_s \sqrt{(n-2)/(1-r_s^2)}$.

To aid practical interpretation, a \textbf{quartile analysis} stratifies the sample into
four approximately equal-sized experience bands and reports the percentage of respondents
in each band who fall into the (a) trust, (b) neutral, and (c) distrust categories.
This non-parametric summary makes any threshold effects visible without distributional
assumptions.

%-----------------------------------------------------------------------
\section{Results}
\label{sec:results}

\subsection{Descriptive Statistics by Trust Group}

Table~\ref{tab:main} presents mean and median professional experience, standard deviation,
and the proportion of respondents with ten or more years of experience across the five
trust groups. A clear monotonic trend is visible: mean experience increases from
$8.75$~years in the \emph{Somewhat trust} group to $10.25$~years in the
\emph{Highly distrust} group, a difference of $1.50$~years. Median experience similarly
rises from $6.0$~years (Somewhat trust, Neither) to $8.0$~years (Highly distrust),
and the proportion of ``senior'' developers ($\geq 10$ years) increases monotonically
from $34.7\%$ to $44.4\%$ across the trust-to-distrust spectrum.

Notably, the \emph{Highly trust} group does not follow the simple monotonic pattern
exactly: its mean ($9.34$~years) is slightly higher than the \emph{Somewhat trust} group
($8.75$~years). This suggests a modest U-shape at the extremes, with very high trust
and very high distrust both associated with somewhat greater experience, while moderate
trust corresponds to the least experienced profile. Despite this nuance, the Spearman
correlation confirms a net positive monotonic relationship between distrust ordinal and
experience (Section~\ref{sec:corr}).

\begin{table}[ht]
\centering
\caption{Professional experience (years) and education across AI accuracy trust levels ($n = 31{,}025$).}
\label{tab:main}
\renewcommand{\arraystretch}{1.15}
\begin{tabular}{lrrrrcc}
\toprule
Trust Level & $n$ & Mean & Median & SD & $\geq$\,10 yrs & Edu (mean) \\
\midrule
Highly trust               &    699 &  9.34 & 6.0 & 9.10 & 38.6\% & 3.59 \\
Somewhat trust             & 12{,}175 & 8.75 & 6.0 & 8.29 & 34.7\% & 3.69 \\
Neither trust nor distrust &  8{,}344 & 8.90 & 6.0 & 8.16 & 36.1\% & 3.73 \\
Somewhat distrust          &  7{,}225 & 9.42 & 7.0 & 8.42 & 39.0\% & 3.64 \\
Highly distrust            &  2{,}582 & 10.25 & 8.0 & 8.44 & 44.4\% & 3.63 \\
\midrule
Total                      & 31{,}025 & 9.12 & 6.0 & 8.33 & 37.1\% & 3.67 \\
\bottomrule
\end{tabular}
\end{table}

The rightmost column of Table~\ref{tab:main} reports the mean education ordinal (0--6
scale) per trust group. Values range narrowly between $3.59$ and $3.73$, corresponding
approximately to the bachelor's degree level throughout. This uniformity confirms that
educational attainment does not differ meaningfully across trust groups and does not
confound the experience--distrust association.

\subsection{Kruskal--Wallis Test}

The Kruskal--Wallis H-test rejects H\textsubscript{0}: $H(4) = 127.78$, $p < 0.001$,
far exceeding the $\chi^2$ critical value of $18.47$ for $\alpha = 0.001$ with
$\text{df} = 4$. The null hypothesis of equal experience distributions across trust
levels is rejected with high confidence. The one-way ANOVA eta-squared estimate is
$\eta^2 = 0.003$, indicating a statistically significant but substantively small effect
--- approximately $0.3\%$ of the variance in experience is attributable to trust group
membership.

The pairwise comparison between \emph{Somewhat trust} ($n = 12{,}175$, mean $= 8.75$~yrs)
and \emph{Highly distrust} ($n = 2{,}582$, mean $= 10.25$~yrs) yields Cohen's $d = 0.18$,
and the comparison between \emph{Highly trust} ($n = 699$, mean $= 9.34$~yrs) and
\emph{Highly distrust} gives $d = 0.11$. Both are classified as small effects under
Cohen's conventions.

\subsection{Spearman Rank Correlation}
\label{sec:corr}

The Spearman correlation between trust ordinal ($1 =$ most trusting, $5 =$ most
distrusting) and years of professional experience is $r_s = 0.102$,
$t(31{,}023) = 18.09$, $p < 0.001$. The corresponding coefficient of determination
is $r_s^2 = 0.010$, indicating that experience explains roughly $1\%$ of the rank
variance in trust orientation. While the effect size is small, the statistical
significance is beyond doubt given the large sample size, and the direction is
consistent: more experience is associated with more distrust.

\subsection{Quartile Analysis}

Table~\ref{tab:quartiles} stratifies the analytic sample into four equal-sized
experience quartiles and reports the distribution of trust, neutrality, and distrust
within each band.

\begin{table}[ht]
\centering
\caption{Trust and distrust rates by professional experience quartile ($n = 31{,}025$).}
\label{tab:quartiles}
\renewcommand{\arraystretch}{1.15}
\begin{tabular}{llrrrr}
\toprule
Quartile & Experience range & $n$ & \% Trust & \% Neutral & \% Distrust \\
\midrule
Q1 & 0--3 years   & 7{,}756 & 53.1\% & 22.7\% & 24.2\% \\
Q2 & 3--6 years   & 7{,}756 & 34.0\% & 32.3\% & 33.7\% \\
Q3 & 6--13 years  & 7{,}756 & 40.7\% & 25.9\% & 33.4\% \\
Q4 & 13--50 years & 7{,}757 & 38.2\% & 26.6\% & 35.2\% \\
\bottomrule
\end{tabular}
\end{table}

The most striking feature of Table~\ref{tab:quartiles} is the contrast between Q1
and the remaining quartiles. Junior developers (0--3 years) exhibit a trust-to-distrust
ratio of $2.2:1$ ($53.1\%$ vs.\ $24.2\%$). By Q2 (3--6 years) this ratio collapses
to near parity ($34.0\%$ vs.\ $33.7\%$, ratio $\approx 1.0:1$). The trust advantage
does not recover in Q3 or Q4; instead, the distrust fraction continues to modestly
increase. The Q1-to-Q2 shift therefore represents the sharpest attitudinal transition
in the dataset: a $19.1$ percentage-point drop in trust and a $9.5$ percentage-point
rise in distrust within approximately the first three years of professional experience.

Beyond Q2, the pattern is more gradual. Q3 (6--13 years) and Q4 ($\geq 13$ years)
differ by only $2.5$ percentage points in distrust, suggesting a threshold effect:
early professional experience rapidly recalibrates trust expectations, while additional
seniority produces more modest marginal shifts.

%-----------------------------------------------------------------------
\section{Discussion and Conclusion}
\label{sec:discussion}

\subsection{Interpretation}

This study demonstrates that professional coding experience is a statistically robust,
if modest, predictor of distrust in AI output accuracy among developers who actively
use AI tools. H\textsubscript{1} is supported: mean experience increases monotonically
(with minor deviation at the \emph{Highly trust} extreme) from the most trusting to
the most distrusting group, and this difference is significant by both the Kruskal--Wallis
test and Spearman correlation.

Several mechanisms may account for this finding. First, experienced developers
accumulate larger internal libraries of correct code patterns, enabling faster and more
accurate identification of AI-generated errors. Vaithilingam et al.\
\cite{vaithilingam2022expectation} provided controlled experimental evidence for this
mechanism; the present study shows its signature in observational population-level data
at a scale three orders of magnitude larger than any single usability study. Second,
senior developers disproportionately work on complex, domain-specific, or legacy
codebases where contemporary LLMs underperform \cite{ziegler2022productivity}. Their
AI tools fail more often in practice, driving down trust. Third, experienced practitioners
have lived through multiple technology hype cycles --- from CASE tools and expert systems
to agile frameworks and cloud containerisation --- and may bring a natural wariness
toward claims of transformative productivity gain \cite{bird2022copilot}.

The threshold effect identified in the quartile analysis (the sharp Q1-to-Q2 transition
at approximately three years of experience) aligns with the developmental trajectory
described by \citet{liao2023transparency}: initial AI adoption is enthusiasm-driven, but
sustained interaction produces a more calibrated assessment of reliability. Three years
may represent the approximate period required to encounter a sufficient variety of AI
failures to update one's prior on its accuracy.

The uniformity of education levels across trust groups (Table~\ref{tab:main}) is itself
informative. If distrust were driven primarily by academic training in computer science
or familiarity with statistical modelling, one would expect higher-degree holders to
cluster in the distrust groups. The absence of such a gradient suggests that formal
education does not substitute for lived professional experience in forming accurate
judgements about AI output quality.

\subsection{Limitations}

Several limitations constrain the interpretive scope of these findings. First, the
analysis is \textbf{cross-sectional}: causal inference is not possible. The experience--distrust
association may reflect reverse causation (skeptical individuals remain in the profession
longer), self-selection (those who adopted AI tools early are disproportionately junior),
or unmeasured confounders such as domain specialisation, industry sector, or the specific
AI tool used.

Second, the \textbf{effect size is small}: $r_s = 0.10$ and $\eta^2 = 0.003$ together
explain approximately 1\% of variance. The statistical significance results from the
large sample rather than from a large signal. Practically, this means that professional
experience alone is a poor predictor of individual trust behaviour --- the distribution
within any experience band is wide (SD $\approx 8$--$9$ years).

Third, the \textbf{survey sample is not probability-based}. Stack Overflow users over-represent
developers in certain countries, industries, and experience levels. Junior developers in
large technology companies may be systematically underrepresented compared to their share
of the global developer population.

Fourth, \textbf{trust} is measured as a single self-report item without validation.
Behavioural measures (e.g., actual code review depth, AI suggestion acceptance rates) might
yield different results and would be more directly policy-relevant.

\subsection{Practical Implications}

Despite the small effect size, the findings carry operational implications. Junior
developers are disproportionately trusting, raising the risk of uncritical acceptance
of incorrect AI suggestions. Organisations deploying AI tools should consider targeted
onboarding focused on AI failure modes for developers with fewer than three years of
experience, peer review pairings that leverage the calibration advantage of senior
practitioners, and acceptance-rate metrics stratified by seniority to monitor
over-reliance in real time.

\subsection{Conclusion}

Professional coding experience is a statistically significant predictor of distrust in
AI output accuracy among software developers who use AI tools. Novice developers (0--3 years)
show a $2.2:1$ trust-to-distrust ratio, while senior developers ($\geq 13$ years) approach
parity. The effect is modest in magnitude but highly robust across statistical tests, and
education level does not confound the pattern. These findings are consistent with an
expertise-induced calibration account: experience provides the domain knowledge necessary
to detect AI errors, converting initial enthusiasm into more measured, accurate appraisal.
Future longitudinal research should track whether distrust translates into improved code
quality and whether targeted training can accelerate trust calibration for novice developers.

%-----------------------------------------------------------------------
\bibliographystyle{plain}
\begin{thebibliography}{9}

\bibitem{stackoverflow2024}
Stack Overflow.
\newblock Developer Survey 2024.
\newblock Stack Overflow, 2024.
\newblock \url{https://survey.stackoverflow.co/2024/}

\bibitem{vaithilingam2022expectation}
P.~Vaithilingam, T.~Zhang, and E.~L. Glassman.
\newblock Expectation vs.\ experience: Evaluating the usability of code
  generation tools powered by large language models.
\newblock In \emph{Extended Abstracts of the 2022 {CHI} Conference on Human
  Factors in Computing Systems}, {CHI\,EA}~'22. {ACM}, 2022.
\newblock \url{https://doi.org/10.1145/3491101.3519665}

\bibitem{ziegler2022productivity}
A.~Ziegler, E.~Kalliamvakou, X.~A.~Li, A.~Rice, D.~Rifkin, S.~Simister,
  G.~Sittampalam, and E.~Aftandilian.
\newblock Productivity assessment of neural code completion.
\newblock In \emph{Proceedings of the 6th {ACM} {SIGPLAN} International
  Symposium on Machine Programming}, {MAPS}~'22. {ACM}, 2022.
\newblock \url{https://doi.org/10.1145/3520312.3534864}

\bibitem{bird2022copilot}
C.~Bird, D.~Ford, T.~Zimmermann, N.~Forsgren, E.~Kalliamvakou, T.~Lowdermilk,
  and I.~Gazit.
\newblock Taking flight with {Copilot}: Early experiences and observations.
\newblock \emph{{ACM} Queue}, 20(6), 2022.
\newblock \url{https://doi.org/10.1145/3582083}

\bibitem{liao2023transparency}
Q.~V. Liao and J.~W. Vaughan.
\newblock {AI} transparency in the age of {LLMs}: A human-centered research
  roadmap.
\newblock \emph{Harvard Data Science Review}, 2023.
\newblock \url{https://doi.org/10.1162/99608f92.8036723b}

\bibitem{cai2019hello}
C.~J. Cai, S.~Winter, D.~Steiner, L.~Wilcox, and M.~Terry.
\newblock ``{H}ello {AI}'': Uncovering the onboarding needs of medical
  practitioners for human-{AI} collaborative decision-making.
\newblock \emph{Proceedings of the {ACM} on Human-Computer Interaction},
  3({CSCW}), 2019.
\newblock \url{https://doi.org/10.1145/3359206}

\end{thebibliography}

\end{document}
